% 自 notes/docs/1151/seminar/W09_1109_預習.md（87808f3）由 tools/md_to_tex.py 於 2026-11-09 轉出一次；此後本檔為正本。
\chapter{與前後幾週的關係及課前問題}
\label{ch:relations}

\section{與 W04 至 W08、W10 的對照}
\label{sec:relations-1}

\begin{longtable}[]{@{}
  >{\raggedright\arraybackslash}p{(\linewidth - 4\tabcolsep) * \real{0.1300}}
  >{\raggedright\arraybackslash}p{(\linewidth - 4\tabcolsep) * \real{0.2400}}
  >{\raggedright\arraybackslash}p{(\linewidth - 4\tabcolsep) * \real{0.6300}}@{}}
\caption{本週三篇與 W04 至 W08、W10 指定文獻的對照。}\tabularnewline
\toprule\noalign{}
\begin{minipage}[b]{\linewidth}\raggedright
文獻
\end{minipage} & \begin{minipage}[b]{\linewidth}\raggedright
相關論點
\end{minipage} & \begin{minipage}[b]{\linewidth}\raggedright
與本週三篇的關係
\end{minipage} \\
\midrule\noalign{}
\endhead
\bottomrule\noalign{}
\endlastfoot
W04 Fallon、Bickel\ix{persons}{bickel@Bickel, Alexander M.} & 反多數決困難；法院在憲法意義不確定時的裁判權威 & Kramer\ix{persons}{kramer@Kramer, Larry D.} 631 頁說 1960 年代的自由派學者覺得「the countermajoritarian\ix{terms}{05@\zhlidxcat{解釋理論與自制}!反多數決困難} difficulty」並不那麼困難；Ferejohn\ix{persons}{ferejohn@Ferejohn, John} 與 Pasquino\ix{persons}{pasquino@Pasquino, Pasquale} 25 頁提到司法審查\ix{terms}{01@\zhlidxcat{司法審查權限之歸屬}!司法審查}的反多數決面向；兩篇都未引 Bickel \\
W05 Kelsen\ix{persons}{kelsen@Kelsen, Hans（凱爾森）}、Schmitt\ix{persons}{schmitt@Schmitt, Carl（施密特）} & 憲法的守護者應為何者；消極立法\ix{terms}{04@\zhlidxcat{法院之功能}!消極立法}者 & Ferejohn 與 Pasquino 21 頁註 3 引 Kelsen 1931 年〈Wer soll der Hüter der Verfassung sein?〉的義大利文譯本，30 頁註 13 說 Kelsen 回應 Schmitt 時由消極立法轉為承認憲法法院的積極立法功能；Alexander\ix{persons}{alexander@Alexander, Larry} 與 Schauer\ix{persons}{schauer@Schauer, Frederick} 460 至 461 頁註 17 引 Kelsen 的基本規範\ix{terms}{03@\zhlidxcat{憲法之效力基礎}!基本規範}，說明憲法的地位須在文本之外決定 \\
W06 Waldron\ix{persons}{waldron@Waldron, Jeremy} & 反對違憲審查的核心論證；1354 頁註 22 界定司法至上\ix{terms}{01@\zhlidxcat{司法審查權限之歸屬}!司法至上} & 三篇都未引 Waldron；Waldron 1350 頁註 11\footcite[1350]{waldron2006} 把 Kramer 的著作與 Tushnet\ix{persons}{tushnet@Tushnet, Mark} 並列為近年攻擊違憲審查的著作；W06 講義稱 Kramer 接受違憲審查而反對司法至上，與 Kramer 622 至 623 頁共和派\ix{terms}{01@\zhlidxcat{司法審查權限之歸屬}!共和派}的人民憲政主義\ix{terms}{01@\zhlidxcat{司法審查權限之歸屬}!人民憲政主義}相合 \\
W08 Tushnet & 非司法審查的表現；定紛止爭\ix{terms}{02@\zhlidxcat{定紛止爭論}!定紛止爭}論的經驗前提 & Tushnet 結論回應 Alexander 與 Schauer 的 Reply\pg{註 209 引 464 頁}；Alexander 與 Schauer 458 頁註 12、472 頁註 48 回應 Tushnet；Kramer 629 頁註 30 引 Tushnet 支持立法機關也能認真審議憲法 \\
W10 Law & 司法權與司法審查的理論 & 待下週預習時補 \\
\end{longtable}

\section{課前可以先想的問題}
\label{sec:relations-2}

\begin{enumerate}
\def\labelenumi{\arabic{enumi}.}
\tightlist
\item
  Alexander\ix{persons}{alexander@Alexander, Larry} 與 Schauer\ix{persons}{schauer@Schauer, Frederick} 457 頁註 11 把權威界定為提供與內容無關\ix{terms}{02@\zhlidxcat{定紛止爭論}!與內容無關之理由}之理由，其成立取決於受拘束者長期以定紛\ix{terms}{02@\zhlidxcat{定紛止爭論}!定紛止爭}為準時所犯的錯誤是否較少。Kramer\ix{persons}{kramer@Kramer, Larry D.} 634 頁稱司法至上\ix{terms}{01@\zhlidxcat{司法審查權限之歸屬}!司法至上}是一種意識形態，作用在把公眾接受的均衡點\ix{terms}{05@\zhlidxcat{解釋理論與自制}!均衡點}移向司法權威。兩篇談的都是其他機關與公眾對最高法院的接受，前者從規範上證成，後者從歷史上描述。Kramer 的敘述若成立，是否動搖 Alexander 與 Schauer 的論證？
\item
  Alexander 與 Schauer 457 頁說，最高法院在 Cooper\ix{cases}{02@美國!01@聯邦最高法院!1958Cooper v. Aaron（1958）@Cooper v. Aaron（1958）} 與 City of Boerne\ix{cases}{02@美國!01@聯邦最高法院!1997City of Boerne v. Flores（1997）@City of Boerne v. Flores（1997）} 宣告自身的至上地位，預設了其他部門應接受該宣告的結論，屬於循環論證。釋字第 185 號\ix{cases}{01@中華民國!01@司法院大法官!0185@釋字第 185 號}解釋自行宣示「其所為之解釋，自有拘束全國各機關及人民之效力」\footcite{jy185}，2022 年施行的憲法訴訟法第 38 條\ix{statutes}{01@中華民國!02@憲法訴訟法!038@第 38 條}第 1 項\footcite[第 38 條第 1 項]{csfa}則以法律明定判決拘束各機關及人民。依兩人 459 至 463 頁先於憲法\ix{terms}{03@\zhlidxcat{憲法之效力基礎}!先於憲法的憲法}的憲法之說，中華民國憲法法庭的解釋權威，其根據在憲法第 78 條\ix{statutes}{01@中華民國!01@中華民國憲法!078@第 78 條}\footcite[第 78 條]{roc-constitution}、釋字第 185 號解釋、憲法訴訟法，或各機關與人民實際的接受？
\item
  Kramer 633 頁說，在司法至上之下，自制\ix{terms}{05@\zhlidxcat{解釋理論與自制}!司法自制}只剩不致引發政治反撲的審慎\ix{terms}{05@\zhlidxcat{解釋理論與自制}!審慎}理由。釋字第 328 號\ix{cases}{01@中華民國!01@司法院大法官!0328@釋字第 328 號}解釋\footcite{jy328}以固有疆域的界定為「重大之政治問題」而不予解釋\ix{terms}{06@\zhlidxcat{審查制度之類型}!不予解釋}。這種不予解釋，屬於 Kramer 所說共和派\ix{terms}{01@\zhlidxcat{司法審查權限之歸屬}!共和派}或 Thayer\ix{persons}{thayer@Thayer, James Bradley} 式以其他機關另有解釋權威為前提的自制，或屬於審慎的權宜？
\item
  Alexander 與 Schauer 482 頁註 78 與 W08 Tushnet\ix{persons}{tushnet@Tushnet, Mark} 的結論都說，最高法院與國會何者較能定紛，需要尚未進行的經驗研究。若要做這種研究，應比較哪些指標：裁判被遵守的程度、見解的穩定、立法的修正、公眾的信任？證據不足時，應由主張司法至上者或主張分權解釋者負擔論證？
\item
  Ferejohn\ix{persons}{ferejohn@Ferejohn, John} 與 Pasquino\ix{persons}{pasquino@Pasquino, Pasquale} 33 頁說 Kelsen 型\ix{terms}{06@\zhlidxcat{審查制度之類型}!Kelsen 型模式}法院傾向以單一聲音發言，很少公開聽審或公布不同意見\ix{terms}{04@\zhlidxcat{法院之功能}!不同意見與協同意見}書。大法官解釋與憲法法庭判決都附公布協同意見書與不同意見書。依兩人的三種類型，中華民國憲法法庭的審議較接近 Kelsen\ix{persons}{kelsen@Kelsen, Hans（凱爾森）} 型或美國最高法院？公布意見書對 Ferejohn 與 Pasquino 所說的合法性負擔\ix{terms}{07@\zhlidxcat{憲法審判之審議}!合法性負擔}有何影響？
\item
  Ferejohn 與 Pasquino 21 頁把 Kelsen 型法院的擴散歸因於後威權\ix{terms}{06@\zhlidxcat{審查制度之類型}!後威權政體}國家對法官或議會的不信任。中華民國的大法官設於 1948 年行憲之時，在 1990 年代民主化期間作成釋字第 261 號\ix{cases}{01@中華民國!01@司法院大法官!0261@釋字第 261 號}解釋\footcite{jy261}等裁判。兩人對後威權憲法審判的分析\pg{31 至 32 頁}，能否適用於中華民國大法官在民主轉型時期的角色？
\end{enumerate}
